\documentclass[11pt,a4paper]{article}
\usepackage[margin=2.5cm]{geometry}
\usepackage{hyperref,xcolor,booktabs,fancyhdr,titlesec,listings}
\usepackage[T1]{fontenc}\usepackage{lmodern}

\definecolor{cyber}{HTML}{00FFD6}
\definecolor{mg}{HTML}{FF66AC}
\hypersetup{colorlinks,linkcolor=cyber,urlcolor=cyber}

\pagestyle{fancy}\fancyhf{}
\fancyhead[L]{\textcolor{cyber}{\textbf{FLLC CYBERDECK}}}
\fancyhead[R]{\textcolor{mg}{\thepage}}
\fancyfoot[C]{\small\textcolor{gray}{Local AI Integration Guide}}

\titleformat{\section}{\Large\bfseries\color{cyber}}{\thesection}{1em}{}
\titleformat{\subsection}{\large\bfseries\color{mg}}{\thesubsection}{1em}{}

\lstset{basicstyle=\ttfamily\small,backgroundcolor=\color{black!5},
        frame=single,rulecolor=\color{gray},breaklines=true}

\title{%
  {\Huge\bfseries\textcolor{cyber}{Local AI Integration Guide}}\\[6pt]
  {\Large Ollama \textbullet\ OpenAI-Compatible \textbullet\ None}\\[4pt]
  {\normalsize Version 2.0 --- \today}
}
\author{FLLC Engineering Division}
\date{}

\begin{document}
\maketitle\thispagestyle{fancy}
\tableofcontents\newpage

\section{Overview}
FLLC CyberDeck supports three AI adapter backends.  The application functions
\textbf{fully without any AI backend configured}.

\section{Adapter: None}
Default mode.  All AI endpoints return passthrough results (truncated input,
generic suggestions).  No external calls are made.

\section{Adapter: Ollama}
\subsection{Setup}
\begin{lstlisting}[language=bash]
# Install Ollama
curl -fsSL https://ollama.ai/install.sh | sh

# Pull a model
ollama pull llama3

# Ollama runs on port 11434 by default
\end{lstlisting}

\subsection{Configuration}
In the Settings page or via API:
\begin{lstlisting}
POST /settings/ai
{
  "backend": "ollama",
  "config": {
    "endpoint": "http://localhost:11434",
    "model": "llama3"
  }
}
\end{lstlisting}

\subsection{Health Check}
The adapter checks \texttt{GET /api/tags} on the Ollama endpoint.
If unavailable, it falls back to the NoneAdapter gracefully.

\subsection{Docker Compose}
To run Ollama as part of the stack:
\begin{lstlisting}[language=yaml]
ollama:
  image: ollama/ollama:latest
  ports: ["11434:11434"]
  volumes: ["ollama_data:/root/.ollama"]
\end{lstlisting}

\section{Adapter: OpenAI-Compatible}
Any local server exposing \texttt{/v1/chat/completions}:
\begin{itemize}
  \item LM Studio
  \item text-generation-webui with OpenAI extension
  \item LocalAI
  \item vLLM with OpenAI server mode
\end{itemize}

\subsection{Configuration}
\begin{lstlisting}
POST /settings/ai
{
  "backend": "openai-compatible-local",
  "config": {
    "endpoint": "http://localhost:8080/v1",
    "model": "local-model",
    "api_key": "none"
  }
}
\end{lstlisting}

\section{AI Capabilities}
\begin{tabular}{lp{9cm}}
\toprule
\textbf{Endpoint} & \textbf{Description} \\
\midrule
\texttt{POST /ai/summarize}      & Summarize analyst notes \\
\texttt{POST /ai/draft-finding}  & Draft professional finding language \\
\texttt{POST /ai/suggest-names}  & Suggest component names \\
\texttt{POST /ai/assist-report}  & Generate executive summary text \\
\texttt{POST /ai/cluster}        & Cluster evidence items \\
\bottomrule
\end{tabular}

\section{Safety Constraints}
AI is \textbf{never} used for:
\begin{itemize}
  \item Payload or exploit generation
  \item Attack planning or evasion
  \item Credential harvesting
  \item Any offensive capability
\end{itemize}

\vfill
\begin{center}
\textcolor{gray}{\small\copyright\ FLLC \the\year.  For internal authorized use only.}
\end{center}
\end{document}
